% GENERATED FILE. Edit canonical lessons, then run npm run generate:books.
% canonical-source-hash: fnv1a64:1ebf816f6818f437
% canonical-lessons: PT-C370-a-menos-que, PT-C370-a-nao-ser-que, PT-C370-senao, PT-C370-eventual, PT-C370-eventualmente

\chapter{Caso contrário}
\label{ch:pt-caso-contrario}
\hlchaptermodality{370}

\begin{chapteropening}
\textbf{By the end of this chapter:} \emph{I can talk about what might stop a plan: unless, otherwise, and possible events.}

Complete the last lesson of chapter 370: I can talk about what might stop a plan: unless, otherwise, and possible events.
\end{chapteropening}

\section[a menos que]{a menos que — unless}

\label{lesson:PT-C370-a-menos-que}



\begin{quote}
Before the new one: say the Portuguese for in case, then the Portuguese for depending on.
\end{quote}

\subsection*{You'll want to know: a menos que}
\textbf{a menos que} — \textquotedblleft{}unless\textquotedblright{}.

\begin{quote}
Vamos a pé, a menos que chova.
\end{quote}

\begin{quote}
\emph{We will walk, unless it rains.}
\end{quote}

Takes the subjunctive.

\begin{cousinweb}
\emph{a menos que} is built from \emph{a}, \emph{menos} \textquotedblleft{}less\textquotedblright{}, inherited from Latin \emph{minus} \textquotedblleft{}less\textquotedblright{} (Lewis \& Short), and \emph{que}.
\end{cousinweb}

\begin{grammarlens}[title={What you've built}]
One more word for telling, explaining and describing at B1.
\end{grammarlens}

\subsection*{Your turn}
\begin{itemize}
\raggedright
  \item \emph{Say it:} \emph{a menos que}
  \item \emph{Say it:} \emph{a menos que}, once more
  \item \emph{Say it:} a sentence of your own with \emph{a menos que}
  \item \emph{Recall:} say the Portuguese for in case, then the Portuguese for depending on, then say \emph{a menos que} again
\end{itemize}

\begin{tcolorbox}[breakable,colback=teal!4,colframe=teal!35!black,title={Before you move on}]
What does a menos que mean? (\textquotedblleft{}Unless\textquotedblright{}.) Say it once more.
\end{tcolorbox}
\section[a não ser que]{a não ser que — except if, unless}

\label{lesson:PT-C370-a-nao-ser-que}



\begin{quote}
Before the new one: say the Portuguese for depending on, then the Portuguese for unless.
\end{quote}

\subsection*{You'll want to know: a não ser que}
\textbf{a não ser que} — \textquotedblleft{}except if, unless\textquotedblright{}.

\begin{quote}
Amanhã não trabalho, a não ser que o chefe me chame.
\end{quote}

\begin{quote}
\emph{Tomorrow I am not working, unless the boss calls me.}
\end{quote}

Takes the subjunctive; it means the same as \emph{a menos que}.

\begin{cousinweb}
\emph{a não ser que} is built from \emph{não} \textquotedblleft{}not\textquotedblright{}, inherited from Latin \emph{nōn} (Lewis \& Short), \emph{ser} \textquotedblleft{}to be\textquotedblright{}, and \emph{que}: \textquotedblleft{}if it is not that\textquotedblright{}.
\end{cousinweb}

\begin{grammarlens}[title={What you've built}]
One more word for telling, explaining and describing at B1.
\end{grammarlens}

\subsection*{Your turn}
\begin{itemize}
\raggedright
  \item \emph{Say it:} \emph{a não ser que}
  \item \emph{Say it:} \emph{a não ser que}, once more
  \item \emph{Say it:} a sentence of your own with \emph{a não ser que}
  \item \emph{Recall:} say the Portuguese for depending on, then the Portuguese for unless, then say \emph{a não ser que} again
\end{itemize}

\begin{tcolorbox}[breakable,colback=teal!4,colframe=teal!35!black,title={Before you move on}]
What does a não ser que mean? (\textquotedblleft{}Except if, unless\textquotedblright{}.) Say it once more.
\end{tcolorbox}
\section[senão]{senão — otherwise, or else; but (rather)}

\label{lesson:PT-C370-senao}



\begin{quote}
Before the new one: say the Portuguese for unless, then the Portuguese for except if.
\end{quote}

\subsection*{You'll want to know: senão}
\textbf{senão} — \textquotedblleft{}otherwise, or else; but (rather)\textquotedblright{}.

\begin{quote}
Temos de sair já, senão perdemos o comboio.
\end{quote}

\begin{quote}
\emph{We have to leave now, or else we miss the train.}
\end{quote}

One word for or else; \emph{se não} in two words means if not: \emph{se não chover}.

\begin{cousinweb}
\emph{senão} is built from \emph{se} \textquotedblleft{}if\textquotedblright{}, from Latin \emph{sī}, and \emph{não} \textquotedblleft{}not\textquotedblright{}, from Latin \emph{nōn} (EtymDB): \textquotedblleft{}if not, then\textquotedblright{}.
\end{cousinweb}

\begin{grammarlens}[title={What you've built}]
One more word for telling, explaining and describing at B1.
\end{grammarlens}

\subsection*{Your turn}
\begin{itemize}
\raggedright
  \item \emph{Say it:} \emph{senão}
  \item \emph{Say it:} \emph{senão}, once more
  \item \emph{Say it:} a sentence of your own with \emph{senão}
  \item \emph{Recall:} say the Portuguese for unless, then the Portuguese for except if, then say \emph{senão} again
\end{itemize}

\begin{tcolorbox}[breakable,colback=teal!4,colframe=teal!35!black,title={Before you move on}]
What does senão mean? (\textquotedblleft{}Otherwise, or else; but (rather)\textquotedblright{}.) Say it once more.
\end{tcolorbox}
\section[eventual]{eventual — possible, potential (that may happen)}

\label{lesson:PT-C370-eventual}



\begin{quote}
Before the new one: say the Portuguese for except if, then the Portuguese for otherwise.
\end{quote}

\subsection*{You'll want to know: eventual}
\textbf{eventual} — \textquotedblleft{}possible, potential (that may happen)\textquotedblright{}.

\begin{quote}
Levo um guarda-chuva para uma eventual chuva.
\end{quote}

\begin{quote}
\emph{I am taking an umbrella in case of rain.}
\end{quote}

A false friend: \emph{eventual} means possible, not final.

\begin{cousinweb}
\emph{eventual} is a learned formation on Latin \emph{ēventus} \textquotedblleft{}an occurrence, an event\textquotedblright{} (EtymDB, Lewis \& Short), from \emph{ēvenīre} \textquotedblleft{}to come out, to happen\textquotedblright{}, on \emph{venīre} \textquotedblleft{}to come\textquotedblright{}.
\end{cousinweb}

\begin{grammarlens}[title={What you've built}]
One more word for telling, explaining and describing at B1.
\end{grammarlens}

\subsection*{Your turn}
\begin{itemize}
\raggedright
  \item \emph{Say it:} \emph{eventual}
  \item \emph{Say it:} \emph{eventual}, once more
  \item \emph{Say it:} a sentence of your own with \emph{eventual}
  \item \emph{Recall:} say the Portuguese for except if, then the Portuguese for otherwise, then say \emph{eventual} again
\end{itemize}

\begin{tcolorbox}[breakable,colback=teal!4,colframe=teal!35!black,title={Before you move on}]
What does eventual mean? (\textquotedblleft{}Possible, potential (that may happen)\textquotedblright{}.) Say it once more.
\end{tcolorbox}
\section[eventualmente]{eventualmente — possibly, if need be}

\label{lesson:PT-C370-eventualmente}



\begin{quote}
Before the new one: say the Portuguese for otherwise, then the Portuguese for possible.
\end{quote}

\subsection*{You'll want to know: eventualmente}
\textbf{eventualmente} — \textquotedblleft{}possibly, if need be\textquotedblright{}.

\begin{quote}
Eventualmente, podemos jantar mais tarde.
\end{quote}

\begin{quote}
\emph{If need be, we can have dinner later.}
\end{quote}

A false friend: eventually is \emph{finalmente} or \emph{por fim}.

\begin{cousinweb}
\emph{eventualmente} is built on \emph{eventual} \textquotedblleft{}possible\textquotedblright{} with the adverb ending \emph{-mente}; \emph{eventual} rests on Latin \emph{ēventus} \textquotedblleft{}an occurrence\textquotedblright{} (Lewis \& Short).
\end{cousinweb}

\begin{grammarlens}[title={What you've built}]
One more word for telling, explaining and describing at B1.
\end{grammarlens}

\subsection*{Your turn}
\begin{itemize}
\raggedright
  \item \emph{Say it:} \emph{eventualmente}
  \item \emph{Say it:} \emph{eventualmente}, once more
  \item \emph{Say it:} a sentence of your own with \emph{eventualmente}
  \item \emph{Recall:} say the Portuguese for otherwise, then the Portuguese for possible, then say \emph{eventualmente} again
\end{itemize}

\begin{tcolorbox}[breakable,colback=teal!4,colframe=teal!35!black,title={Before you move on}]
What does eventualmente mean? (\textquotedblleft{}Possibly, if need be\textquotedblright{}.) Say it once more.
\end{tcolorbox}
